% from ega3-1-fr.tex
\oldpage[0\textsubscript{III}]{58}
et il résulte aussitôt de la définition du cup-produit (G, II, 6.6) que ces
homomorphismes commutent au passage de $V$ à un sous-espace ouvert $V'$ de
$V$. D'autre part, on a un homomorphisme d'anneaux
\[
\Gamma(V,\mathcal O_Y)\longrightarrow\Gamma(V,\psi_*(\mathcal O_X))
=\Gamma(f^{-1}(V),\mathcal O_X)
\]
provenant de $\theta$, d'où un homomorphisme canonique de produits tensoriels
\[
\begin{aligned}
&H^p(f^{-1}(V),\mathcal F)\otimes_{\Gamma(V,\mathcal O_Y)}
H^q(f^{-1}(V),\mathcal G)\\
&\qquad\longrightarrow
H^p(f^{-1}(V),\mathcal F)
\otimes_{\Gamma(f^{-1}(V),\mathcal O_X)}
H^q(f^{-1}(V),\mathcal G)
\end{aligned}
\]
qui est lui aussi compatible avec la restriction de $V'$ à $V$. Par
composition, on obtient donc un homomorphisme de
$\Gamma(V,\mathcal O_Y)$-modules, qui définit un homomorphisme canonique
fonctoriel en $\mathcal F$ et $\mathcal G$ pour les faisceaux associés aux
préfaisceaux considérés :
\[
\label{0.12.2.2.1-fr}
\tag{12.2.2.1}
R^p f_*(\mathcal F)\otimes_{\mathcal O_Y}R^q f_*(\mathcal G)
\longrightarrow
R^{p+q}f_*(\mathcal F\otimes_{\mathcal O_X}\mathcal G).
\]
On notera que pour $p=q=0$, cet homomorphisme se réduit à
(\hyperref[0.4.2.2.1-fr]{0\textsubscript{I}, 4.2.2.1}).
\end{env}

\begin{proposition}[12.2.3]
\label{0.12.2.3-fr}
\emph{Pour tout $\mathcal O_X$-Module $\mathcal F$ et tout
$\mathcal O_Y$-Module localement libre de rang fini $\mathcal L$, on a des
isomorphismes canoniques fonctoriels
\[
\label{0.12.2.3.1-fr}
\tag{12.2.3.1}
R^p f_*(\mathcal F)\otimes_{\mathcal O_Y}\mathcal L
\xrightarrow{\sim}
R^p f_*(\mathcal F\otimes_{\mathcal O_X}f^*(\mathcal L)).
\]}
\end{proposition}

\begin{proof}
L'homomorphisme (\hyperref[0.12.2.3.1-fr]{12.2.3.1}) s'obtient en composant
l'homomorphisme, cas particulier de
(\hyperref[0.12.2.2.1-fr]{12.2.2.1}) :
\[
\label{0.12.2.3.2-fr}
\tag{12.2.3.2}
R^p f_*(\mathcal F)\otimes_{\mathcal O_Y}f_*(f^*(\mathcal L))
\longrightarrow
R^p f_*(\mathcal F\otimes_{\mathcal O_X}f^*(\mathcal L))
\]
avec l'homomorphisme du premier membre de
(\hyperref[0.12.2.3.1-fr]{12.2.3.1}) dans celui de
(\hyperref[0.12.2.3.2-fr]{12.2.3.2}), provenant de l'homomorphisme canonique
(\hyperref[0.4.4.3.2-fr]{0\textsubscript{I}, 4.4.3.2}). Pour vérifier que
(\hyperref[0.12.2.3.1-fr]{12.2.3.1}) est un isomorphisme lorsque $\mathcal L$
est localement libre, on peut aussitôt se ramener au cas
$\mathcal L=\mathcal O_Y$, la question étant locale sur $Y$, et les foncteurs
envisagés étant additifs en $\mathcal L$. Mais alors, la proposition se
ramène, vu la définition de (\hyperref[0.12.2.2.1-fr]{12.2.2.1}), à la
vérification du fait que l'homomorphisme correspondant de préfaisceaux est
bijectif, ce qui est immédiat en vertu de la relation
$f^*(\mathcal O_Y)=\mathcal O_X$.
\end{proof}

\begin{env}[12.2.4]
\label{0.12.2.4-fr}
Soient $(Z,\mathcal O_Z)$ un troisième espace annelé, $g:Y\to Z$ un
morphisme d'espaces annelés. On sait (G, II, 7.1 et 3.1.1) que pour tout
$\mathcal O_X$-Module injectif $\mathcal G$, $f_*(\mathcal G)$ est un
faisceau flasque de groupes abéliens, et par suite
(\hyperref[0.12.2.1-fr]{12.2.1}) on a
$R^p g_*(f_*(\mathcal G))=0$ pour tout $p>0$. Il s'ensuit (T, 2.4.1) que la
suite spectrale de Leray des foncteurs composés est applicable au foncteur
composé $g_*f_*$ : il y a une suite spectrale birégulière dont
l'aboutissement est le foncteur $R^\bullet h_*$, où $h=g\circ f$, et dont le
terme $E_2$ est donné par
\[
\label{0.12.2.4.1-fr}
\tag{12.2.4.1}
E_2^{pq}=R^p g_*(R^q f_*(\mathcal F)).
\]
\end{env}

\begin{env}[12.2.5]
\label{0.12.2.5-fr}
Sous les conditions de (\hyperref[0.12.2.4-fr]{12.2.4}), nous allons définir
directement des homomorphismes canoniques de $\mathcal O_Z$-Modules
\[
\label{0.12.2.5.1-fr}
\tag{12.2.5.1}
R^n g_*(f_*(\mathcal F))\longrightarrow R^n h_*(\mathcal F)
\]
\[
\label{0.12.2.5.2-fr}
\tag{12.2.5.2}
R^n h_*(\mathcal F)\longrightarrow g_*(R^n f_*(\mathcal F))
\]
